\documentclass[11pt]{article}
\usepackage{ap-psych}
\title{Unit 2: Cognition}
\author{Paul Wang}
\date{November 2026}
\docpart{Unit 2 \textperiodcentered\ Study Guide \textperiodcentered\ 15--25\% of MCQ}

\begin{document}
\maketitle

\begin{abstract}
Unit 2 follows information from the moment it is perceived to the moment it is used or lost. It covers perception and attention (2.1); thinking, problem-solving, and decision-making biases (2.2); how memory is organized (2.3); encoding (2.4), storage (2.5), retrieval (2.6), and forgetting and distortion (2.7); and intelligence and its measurement (2.8). This is the unit with the most look-alike terms, so the traps flag the pairs that are tested against each other. Each section ends with check questions; the key is on the last page.
\end{abstract}

\scoringline{Check questions: \blank\ / \qtotal \qquad Target: 19+ \qquad Pairs I mix up: \blank[10em]}

\section{Perception (2.1)}

\subsection{Processing and expectations}
\textbf{Bottom-up processing} builds a perception from raw sensory data. \textbf{Top-down processing} uses expectations, knowledge, and context. \textbf{Schemas} are mental frameworks that organize information; a \textbf{perceptual set} is a readiness to perceive something in a particular way. Context, motivation, emotion, and culture all shape perception.

\subsection{Gestalt principles}
The whole differs from the sum of its parts. \textbf{Figure--ground} (separating an object from its background), \textbf{proximity} (near items grouped), \textbf{similarity} (alike items grouped), \textbf{closure} (filling gaps to see a complete shape).

\subsection{Attention}
\textbf{Selective attention}: focusing on one stimulus; the \textbf{cocktail party effect} is hearing your name across a noisy room. \textbf{Inattentional blindness}: failing to see an obvious object when attention is elsewhere (the unnoticed gorilla). \textbf{Change blindness}: failing to notice a change in the scene.

\subsection{Depth, constancy, and motion}
\begin{compare}{Cue type}{Cue}{How it works}
Binocular & Retinal disparity & The two eyes' images differ more for nearer objects. \\
 & Convergence & Eyes turn inward more for near objects (muscle strain). \\
Monocular & Relative size & Smaller image of a same-size object = farther. \\
 & Interposition & An object that blocks another is closer. \\
 & Relative clarity & Hazy objects look farther. \\
 & Texture gradient & Texture becomes finer and denser with distance. \\
 & Linear perspective & Parallel lines converge in the distance. \\
\end{compare}
\textbf{Perceptual constancy}: perceiving objects as stable in \textbf{size}, \textbf{shape}, and \textbf{color/lightness} despite changing images. \textbf{Apparent motion}: perceiving movement in a rapid series of still images (the phi phenomenon, flip-books, film).

\begin{q}{2}{2.1 Top-down}{Expectations from prior knowledge shaping perception = top-down processing.}
A proofreader misses the typo in ``the the'' because she expects normal English. This is an example of
\chh{bottom-up processing}{top-down processing}{change blindness}{convergence}
\end{q}

\begin{q}{4}{2.1 Attention}{Missing a visible, unexpected object while focused elsewhere = inattentional blindness.}
While counting basketball passes, many viewers fail to notice a person in a gorilla suit walking through the scene. This demonstrates
\ch{change blindness}{the cocktail party effect}{perceptual constancy}{inattentional blindness}
\end{q}

\begin{q}{1}{2.1 Depth cues}{Parallel lines converging with distance = linear perspective, a monocular cue.}
A painter makes railroad tracks meet at the horizon to show distance. She is using
\chh{linear perspective}{retinal disparity}{convergence}{closure}
\end{q}

\begin{q}{3}{2.1 Constancy}{Perceiving a door as rectangular as its image becomes a trapezoid = shape constancy.}
As a door swings open, its image on your retina becomes a trapezoid, but you still see it as rectangular. This is
\chh{size constancy}{apparent motion}{shape constancy}{figure--ground}
\end{q}

\section{Thinking, problem-solving, and decisions (2.2)}

\textbf{Concepts} group similar things; a \textbf{prototype} is the best example (a robin for ``bird''). An \textbf{algorithm} is a step-by-step procedure that guarantees a solution; a \textbf{heuristic} is a mental shortcut that is fast but error-prone. \textbf{Executive functions} (prefrontal cortex) plan, organize, and inhibit.

\begin{terms}
Representativeness heuristic & Judging likelihood by how well something matches a prototype, ignoring base rates (``She's quiet and reads, so she's a librarian, not a salesperson''). \\
Availability heuristic & Judging likelihood by how easily examples come to mind (fearing plane crashes after news coverage). \\
Mental set & Sticking with a strategy that worked before even when a better one exists. \\
Functional fixedness & Seeing objects only in their usual function (not using a coin as a screwdriver). \\
Priming & Earlier exposure to a stimulus unconsciously influences a later response. \\
Framing & How an issue is worded changes decisions (``90\% survival'' vs.\ ``10\% mortality''). \\
Gambler's fallacy & Believing past random events change future odds (``Red is due''). \\
Sunk-cost fallacy & Continuing an investment because of what has already been spent. \\
Confirmation bias & Seeking information that supports existing beliefs, ignoring the rest. \\
\end{terms}

\textbf{Creativity} draws on \textbf{divergent thinking} (generating many possible answers) more than \textbf{convergent thinking} (narrowing to the single best answer).

\begin{trap}
Representativeness vs.\ availability: if the judgment rests on \emph{resemblance to a type}, it is representativeness; if it rests on \emph{how easily examples come to mind}, it is availability.
\end{trap}

\begin{q}{2}{2.2 Heuristics}{Judging risk by vivid, easily recalled news = availability heuristic.}
After watching several news stories about shark attacks, Owen believes shark attacks are more common than car crashes. He is using
\ch{the representativeness heuristic}{the availability heuristic}{an algorithm}{functional fixedness}
\end{q}

\begin{q}{4}{2.2 Biases}{Continuing because of money already spent = sunk-cost fallacy.}
Leila keeps going to a gym class she hates because she already paid for the full year. This reflects
\chh{framing}{mental set}{the gambler's fallacy}{the sunk-cost fallacy}
\end{q}

\begin{q}{1}{2.2 Problem-solving}{Seeing an object only in its normal use blocks the solution: functional fixedness.}
Needing to hang a picture, Sam does not think to use his shoe as a hammer. This is
\chh{functional fixedness}{confirmation bias}{priming}{divergent thinking}
\end{q}

\begin{q}{3}{2.2 Creativity}{Generating many uses for an object measures divergent thinking.}
``List as many uses for a paperclip as you can in two minutes'' measures
\ch{convergent thinking}{an algorithm}{divergent thinking}{executive function deficits}
\end{q}

\section{How memory is organized (2.3)}

\begin{terms}
Explicit (declarative) & Memories you can consciously recall and state. \textbf{Episodic}: personal events (``my 16th birthday''). \textbf{Semantic}: facts and general knowledge (``Paris is in France''). \\
Implicit (nondeclarative) & Memories retained without conscious recall. \textbf{Procedural}: skills (riding a bike); also classically conditioned responses and priming. \\
Prospective memory & Remembering to do something in the future (take medicine at 8). \\
Long-term potentiation & Repeated stimulation strengthens a synapse; the neural basis of learning and memory. \\
\end{terms}

\textbf{Multi-store model} (Atkinson--Shiffrin): \textbf{sensory memory} \(\rightarrow\) (attention) \(\rightarrow\) \textbf{short-term memory} \(\rightarrow\) (encoding/rehearsal) \(\rightarrow\) \textbf{long-term memory}.
\textbf{Working memory model} (Baddeley): short-term memory is an active workspace with a \textbf{central executive} (directs attention), a \textbf{phonological loop} (sounds and words), and a \textbf{visuospatial sketchpad} (images and locations).
\textbf{Levels of processing}: deeper (meaning-based) processing produces more durable memories than shallow (appearance or sound) processing.

\begin{q}{2}{2.3 Memory types}{Fact knowledge not tied to a personal event = semantic memory.}
Knowing that the mitochondrion is the powerhouse of the cell is an example of
\chh{episodic memory}{semantic memory}{procedural memory}{prospective memory}
\end{q}

\begin{q}{3}{2.3 Implicit memory}{H.M.-type patients learn skills (procedural, implicit) without remembering practicing.}
A patient with severe amnesia improves at mirror-tracing each day but never remembers practicing. His improvement reflects
\chh{episodic memory}{semantic memory}{implicit procedural memory}{prospective memory}
\end{q}

\section{Encoding (2.4)}

\textbf{Automatic processing} encodes information without effort (where you parked, how often you saw a word). \textbf{Effortful processing} requires attention and practice.
\begin{itemize}
\item \textbf{Shallow} processing: structural (how it looks) or phonemic (how it sounds). \textbf{Deep} processing: semantic (what it means). The \textbf{self-reference effect}: relating material to yourself boosts recall.
\item \textbf{Mnemonics}: memory aids, including the \textbf{method of loci} (placing items along a familiar route) and acronyms.
\item \textbf{Chunking}: grouping items into meaningful units (FBI-CIA-NBA). \textbf{Hierarchies} organize concepts into categories.
\item \textbf{Spacing effect}: distributed practice beats massed practice (cramming) for long-term retention.
\item \textbf{Serial position effect}: better recall of the first items (\textbf{primacy}, already in long-term memory) and last items (\textbf{recency}, still in short-term memory).
\end{itemize}

\begin{q}{1}{2.4 Encoding}{Thinking about meaning = deep (semantic) processing, best for recall.}
Which study question would produce the best recall of a vocabulary list?
\ch{``Does this word describe you?''}{``Is this word in capital letters?''}{``Does this word rhyme with \emph{train}?''}{``How many letters does this word have?''}
\end{q}

\begin{q}{4}{2.4 Serial position}{Immediate recall favors first (primacy) and last (recency) items.}
Immediately after hearing a 15-word list, participants recall the first few and last few words best. This is
\ch{the spacing effect}{chunking}{proactive interference}{the serial position effect}
\end{q}

\begin{q}{2}{2.4 Spacing}{Distributed practice beats massed practice for long-term retention.}
Which schedule will most likely produce the best exam performance a month later?
\ch{Five hours of study the night before}{One hour a day for five days}{Rereading notes during the exam week only}{Highlighting the textbook once}
\end{q}

\section{Storage (2.5)}

\begin{compare}{Store}{Capacity and duration}{Notes}
Sensory & Large; under 1 s (iconic, visual) to a few s (echoic, auditory). & Sperling's partial-report studies showed brief, complete iconic images. \\
Short-term / working & About 7 \(\pm\) 2 items (fewer for complex items); about 20 s without rehearsal. & \textbf{Maintenance rehearsal} keeps items active. \\
Long-term & Essentially unlimited; potentially permanent. & Memories are stored across networks; \textbf{consolidation} (much of it during sleep) stabilizes them. \\
\end{compare}
Explicit memories depend on the \textbf{hippocampus} (then the cortex); implicit procedural memories depend on the \textbf{cerebellum} and basal ganglia; emotional memories involve the \textbf{amygdala}. \textbf{Highly superior autobiographical memory} lets a rare few recall nearly every day of their lives.

\begin{q}{3}{2.5 Short-term}{STM holds about 7 ± 2 items for about 20 s without rehearsal.}
Without repeating it, you will most likely lose a new phone number from short-term memory within about
\chh{1 second}{5 minutes}{20 seconds}{1 day}
\end{q}

\section{Retrieval (2.6)}

\textbf{Recall} retrieves information without cues (fill-in-the-blank); \textbf{recognition} identifies previously learned items (multiple choice) and is easier. \textbf{Retrieval cues} help; \textbf{encoding specificity} says cues present at encoding help most at retrieval.
\begin{terms}
Context-dependent memory & Recall is better in the same physical setting as learning. \\
State-dependent memory & Recall is better in the same internal state (mood, sobriety). \\
Mood-congruent memory & Current mood cues memories with the same emotional tone. \\
Testing effect & Practicing retrieval strengthens memory more than rereading. \\
Metacognition & Thinking about your own thinking, e.g., judging how well you know material. \\
\end{terms}

\begin{q}{1}{2.6 Retrieval}{Being in the same place as learning aids recall: context-dependent memory.}
A diver recalls a list learned underwater better when tested underwater than on the beach. This illustrates
\ch{context-dependent memory}{state-dependent memory}{mood-congruent memory}{the testing effect}
\end{q}

\begin{q}{4}{2.6 Retrieval}{Current mood cues same-valence memories: mood-congruent memory.}
When she feels sad, Tara easily remembers other sad times in her life. This is
\chh{state-dependent memory}{retroactive interference}{source amnesia}{mood-congruent memory}
\end{q}

\section{Forgetting and distortion (2.7)}

The \textbf{forgetting curve} (Ebbinghaus): most forgetting happens soon after learning, then levels off. Forgetting can be \textbf{encoding failure} (never stored), \textbf{storage decay}, or \textbf{retrieval failure} (the \textbf{tip-of-the-tongue} phenomenon).
\begin{terms}
Proactive interference & Old learning disrupts recall of new information (\emph{pro} = forward). \\
Retroactive interference & New learning disrupts recall of old information (\emph{retro} = backward). \\
Anterograde amnesia & Cannot form \emph{new} memories after an injury. \\
Retrograde amnesia & Cannot recall memories from \emph{before} an injury. \\
Infantile amnesia & Few memories from before about age 3--4. \\
Alzheimer's disease & Progressive memory loss linked to loss of ACh-producing neurons and brain tissue. \\
Repression & In psychodynamic theory, pushing anxiety-provoking memories out of awareness (controversial). \\
Misinformation effect & Misleading information after an event alters the memory (Loftus: ``smashed'' vs.\ ``hit''). \\
Imagination inflation & Imagining an event makes people more confident it happened. \\
Source amnesia & Remembering information but not where it came from. \\
Constructive memory & Memories are rebuilt at each retrieval, not replayed like video. \\
\end{terms}

\begin{trap}
Interference direction: ask \emph{which} memory is being lost. If the new one is lost, the old one interfered forward: proactive. If the old one is lost, the new one interfered backward: retroactive.
\end{trap}

\begin{q}{2}{2.7 Interference}{Old password blocks the new one = proactive interference.}
Mateo keeps typing his old password after changing it. This is
\ch{retroactive interference}{proactive interference}{anterograde amnesia}{encoding failure}
\end{q}

\begin{q}{3}{2.7 Misinformation}{Leading wording after the event altered the memory: misinformation effect.}
Witnesses asked how fast cars were going when they ``smashed'' later report seeing broken glass that was not there. This is the
\chh{forgetting curve}{testing effect}{misinformation effect}{serial position effect}
\end{q}

\begin{q}{1}{2.7 Amnesia}{Losing memories from before the injury = retrograde amnesia.}
After a concussion, a quarterback cannot remember the plays from the hour before the hit but remembers everything since. This is
\chh{retrograde amnesia}{anterograde amnesia}{infantile amnesia}{repression}
\end{q}

\section{Intelligence and achievement (2.8)}

\begin{terms}
General intelligence (\(g\)) & Spearman: a single factor underlies performance on all mental tests. \\
Multiple intelligences & Gardner: several relatively independent abilities (e.g., linguistic, spatial, interpersonal). Popular, weak evidence. \\
Fluid vs.\ crystallized & \textbf{Fluid}: reasoning speed and novel problem-solving; declines with age. \textbf{Crystallized}: accumulated knowledge and vocabulary; grows with age. \\
Intelligence quotient & Original: (mental age \(\div\) chronological age) \(\times 100\). Modern tests (Stanford-Binet, \textbf{Wechsler}/WAIS) use deviation scores, mean 100, SD 15. \\
Achievement vs.\ aptitude & Achievement tests measure what you have learned (AP exams); aptitude tests predict future performance. \\
Standardization & Uniform procedures and norms from a representative pretested group. \\
Reliability & Consistency: test--retest, split-half. \\
Validity & Measures what it claims: \textbf{construct} validity (measures the intended concept), \textbf{predictive} validity (predicts the future outcome it should). \\
Flynn effect & Average IQ scores rose over the 20th century, too fast for genetic change. \\
Stereotype threat & Worry about confirming a negative group stereotype lowers performance. \\
Growth mindset & Believing abilities can grow with effort (Dweck), vs.\ a fixed mindset. \\
\end{terms}
Intelligence tests have been misused (e.g., by eugenicists); culturally biased content or norms can produce misleading group differences.

\begin{q}{4}{2.8 Psychometrics}{Consistent scores on retesting = test-retest reliability (not validity).}
A student takes the same personality test twice, a month apart, and gets nearly identical scores. This shows the test has
\chh{construct validity}{predictive validity}{standardization}{test--retest reliability}
\end{q}

\begin{q}{2}{2.8 Intelligence}{Accumulated knowledge like vocabulary = crystallized intelligence, which holds or grows with age.}
A 70-year-old outperforms her grandson on a vocabulary test but is slower at solving new puzzles. Her vocabulary reflects
\chh{fluid intelligence}{crystallized intelligence}{aptitude}{the Flynn effect}
\end{q}

\begin{q}{3}{2.8 Stereotype threat}{Reminding people of a negative stereotype before a test lowers performance.}
Women told that a math test ``shows gender differences'' score lower than women told it does not. This demonstrates
\chh{the Flynn effect}{a growth mindset}{stereotype threat}{predictive validity}
\end{q}

\makekey
\end{document}
